\ssr{ЗАКЛЮЧЕНИЕ}

Цель работы достигнута: выполнен сравнительный анализ алгоритмов умножения матриц~--- стандартного алгоритма, алгоритма Винограда и трёх его оптимизированных версий~--- по трудоёмкости и по данным замеров процессорного времени.

Все задачи решены:
\begin{enumerate}
	\item описана решаемая согласно варианту задача умножения матриц;
	\item разработаны стандартный алгоритм умножения матриц, алгоритм Винограда и три оптимизированные по варианту версии алгоритма Винограда: с заменой деления на 2 побитовым сдвигом (оптимизация А), с использованием составного оператора \texttt{+=} (оптимизация Б) и с обеими оптимизациями одновременно;
	\item разработанные алгоритмы описаны схемами алгоритмов;
	\item описаны требования к программному обеспечению;
	\item описаны средства реализации алгоритмов, в качестве языка программирования выбран Swift;
	\item разработанные алгоритмы реализованы;
	\item выполнено тестирование реализаций алгоритмов, все реализации прошли функциональные тесты;
	\item описана модель вычислений и выполнен теоретический расчёт трудоёмкости всех алгоритмов;
	\item описаны инструменты замера процессорного времени: функция \texttt{clock\_gettime\_nsec\_np} с часами \texttt{CLOCK\_THREAD\_CPUTIME\_ID};
	\item выполнены массированные замеры процессорного времени для квадратных матриц чётных размеров от 10 до 200 и нечётных размеров от 11 до 201;
	\item выполнен сравнительный анализ алгоритмов по трудоёмкости и по данным замеров.
\end{enumerate}

По результатам сравнительного анализа можно сделать следующие выводы:
\begin{itemize}
	\item все рассмотренные алгоритмы имеют кубическую трудоёмкость, что подтверждается замерами;
	\item по данным замеров для сборки без оптимизаций компилятора все версии алгоритма Винограда работают быстрее стандартного алгоритма: базовая версия~--- в $1{,}6$ раза, версии с оптимизацией Б и с оптимизациями А и Б~--- в $1{,}7$--$1{,}8$ раза при $N = 200$;
	\item преимущество алгоритма Винограда объясняется вдвое меньшим числом обращений к накапливаемому элементу результирующей матрицы; модель вычислений, в которой обращение к элементу массива имеет такую же трудоёмкость, как сложение, это преимущество не предсказывает ($15N^3$ против $14N^3$ в старшем члене);
	\item оптимизация А (замена деления на 2 сдвигом) не даёт измеримого выигрыша, так как граница цикла в Swift вычисляется один раз при входе в цикл;
	\item оптимизация Б (составной оператор \texttt{+=}) уменьшает время работы алгоритма Винограда на 11\,\%, что близко к предсказанным моделью 10\,\%;
	\item совместное применение оптимизаций А и Б даёт такой же выигрыш, как одна оптимизация Б;
	\item различий в соотношении алгоритмов для чётных и нечётных размеров матриц не обнаружено.
\end{itemize}
